% GENERATED FILE. Edit canonical lessons, then run npm run generate:books.
% canonical-source-hash: fnv1a64:cdd663ef34d8ab51
% canonical-lessons: BN-C91-agamikal, BN-C91-sombar, BN-C91-budhbar, BN-C91-brihospotibar, BN-C91-shukrobar

\chapter{Tomorrow, Monday, Wednesday, Thursday, Friday}
\label{ch:bn-tomorrow-monday-wednesday-thursday-friday}
\hlchaptermodality{91}

\begin{chapteropening}
\textbf{By the end of this chapter:} \emph{I can say tomorrow, Monday, Wednesday, Thursday and Friday.}

Complete the last lesson of chapter 91: I can say tomorrow, Monday, Wednesday, Thursday and Friday.
\end{chapteropening}

\section[\texorpdfstring{\bn{আগামীকাল} (āgāmikāl)}{আগামীকাল (āgāmikāl)}]{\bn{আগামীকাল} (āgāmikāl) — tomorrow}

\label{lesson:BN-C91-agamikal}



\begin{quote}
Before the new one: say the Bengali for the day after tomorrow, then the Bengali for yesterday.
\end{quote}

\subsection*{You'll want to know: \bn{আগামীকাল}}
\textbf{\bn{আগামীকাল}} — \emph{āgāmikāl} — \textquotedblleft{}tomorrow\textquotedblright{}.

\begin{cousinweb}
\textbf{\bn{আগামীকাল}} (\emph{āgāmikāl}) is \emph{āgāmi}, \textquotedblleft{}coming\textquotedblright{}, Sanskrit \emph{āgāmin} taken over whole, + \emph{kāl}, \textquotedblleft{}yesterday, tomorrow\textquotedblright{}, from Sanskrit \emph{kalya}, \textquotedblleft{}dawn\textquotedblright{} (CDIAL 3104).
\end{cousinweb}

\begin{grammarlens}[title={What you've built}]
One more word for this chapter.
\end{grammarlens}

\subsection*{Your turn}
\begin{itemize}
\raggedright
  \item \emph{Say it:} \emph{āgāmikāl}
  \item \emph{Say it:} \emph{āgāmikāl}, once more
  \item \emph{Say it:} say \emph{āgāmikāl}
  \item \emph{Recall:} say the Bengali for the day after tomorrow, then the Bengali for yesterday, then say \emph{āgāmikāl} again
\end{itemize}

\begin{tcolorbox}[breakable,colback=teal!4,colframe=teal!35!black,title={Before you move on}]
What does \bn{আগামীকাল} mean? (\textquotedblleft{}Tomorrow\textquotedblright{}.) Say it once more.
\end{tcolorbox}
\section[\texorpdfstring{\bn{সোমবার} (sombār)}{সোমবার (sombār)}]{\bn{সোমবার} (sombār) — Monday}

\label{lesson:BN-C91-sombar}



\begin{quote}
Before the new one: say the Bengali for yesterday, then the Bengali for tomorrow.
\end{quote}

\subsection*{You'll want to know: \bn{সোমবার}}
\textbf{\bn{সোমবার}} — \emph{sombār} — \textquotedblleft{}Monday\textquotedblright{}.

\begin{cousinweb}
\textbf{\bn{সোমবার}} (\emph{sombār}) is Sanskrit \emph{somavāra}, \textquotedblleft{}Monday\textquotedblright{}, taken over whole as a learned loan, a tatsama. It is \emph{soma}, \textquotedblleft{}the moon\textquotedblright{}, + \emph{vāra}, \textquotedblleft{}day of the week\textquotedblright{}.
\end{cousinweb}

\begin{grammarlens}[title={What you've built}]
One more word for this chapter.
\end{grammarlens}

\subsection*{Your turn}
\begin{itemize}
\raggedright
  \item \emph{Say it:} \emph{sombār}
  \item \emph{Say it:} \emph{sombār}, once more
  \item \emph{Say it:} say \emph{sombār}
  \item \emph{Recall:} say the Bengali for yesterday, then the Bengali for tomorrow, then say \emph{sombār} again
\end{itemize}

\begin{tcolorbox}[breakable,colback=teal!4,colframe=teal!35!black,title={Before you move on}]
What does \bn{সোমবার} mean? (\textquotedblleft{}Monday\textquotedblright{}.) Say it once more.
\end{tcolorbox}
\section[\texorpdfstring{\bn{বুধবার} (budhbār)}{বুধবার (budhbār)}]{\bn{বুধবার} (budhbār) — Wednesday}

\label{lesson:BN-C91-budhbar}



\begin{quote}
Before the new one: say the Bengali for tomorrow, then the Bengali for Monday.
\end{quote}

\subsection*{You'll want to know: \bn{বুধবার}}
\textbf{\bn{বুধবার}} — \emph{budhbār} — \textquotedblleft{}Wednesday\textquotedblright{}.

\begin{cousinweb}
\textbf{\bn{বুধবার}} (\emph{budhbār}) is Sanskrit \emph{budhavāra}, \textquotedblleft{}Wednesday\textquotedblright{}, taken over whole as a learned loan, a tatsama. It is \emph{budha}, \textquotedblleft{}Mercury\textquotedblright{}, + \emph{vāra}, \textquotedblleft{}day of the week\textquotedblright{}.
\end{cousinweb}

\begin{grammarlens}[title={What you've built}]
One more word for this chapter.
\end{grammarlens}

\subsection*{Your turn}
\begin{itemize}
\raggedright
  \item \emph{Say it:} \emph{budhbār}
  \item \emph{Say it:} \emph{budhbār}, once more
  \item \emph{Say it:} say \emph{budhbār}
  \item \emph{Recall:} say the Bengali for tomorrow, then the Bengali for Monday, then say \emph{budhbār} again
\end{itemize}

\begin{tcolorbox}[breakable,colback=teal!4,colframe=teal!35!black,title={Before you move on}]
What does \bn{বুধবার} mean? (\textquotedblleft{}Wednesday\textquotedblright{}.) Say it once more.
\end{tcolorbox}
\section[\texorpdfstring{\bn{বৃহস্পতিবার} (brihôspôtibār)}{বৃহস্পতিবার (brihôspôtibār)}]{\bn{বৃহস্পতিবার} (brihôspôtibār) — Thursday}

\label{lesson:BN-C91-brihospotibar}



\begin{quote}
Before the new one: say the Bengali for Monday, then the Bengali for Wednesday.
\end{quote}

\subsection*{You'll want to know: \bn{বৃহস্পতিবার}}
\textbf{\bn{বৃহস্পতিবার}} — \emph{brihôspôtibār} — \textquotedblleft{}Thursday\textquotedblright{}.

\begin{cousinweb}
\textbf{\bn{বৃহস্পতিবার}} (\emph{brihôspôtibār}) is Sanskrit \emph{bṛhaspativāra}, \textquotedblleft{}Thursday\textquotedblright{}, taken over whole as a learned loan, a tatsama. It is \emph{bṛhaspati}, \textquotedblleft{}Jupiter\textquotedblright{}, + \emph{vāra}, \textquotedblleft{}day of the week\textquotedblright{}.
\end{cousinweb}

\begin{grammarlens}[title={What you've built}]
One more word for this chapter.
\end{grammarlens}

\subsection*{Your turn}
\begin{itemize}
\raggedright
  \item \emph{Say it:} \emph{brihôspôtibār}
  \item \emph{Say it:} \emph{brihôspôtibār}, once more
  \item \emph{Say it:} say \emph{brihôspôtibār}
  \item \emph{Recall:} say the Bengali for Monday, then the Bengali for Wednesday, then say \emph{brihôspôtibār} again
\end{itemize}

\begin{tcolorbox}[breakable,colback=teal!4,colframe=teal!35!black,title={Before you move on}]
What does \bn{বৃহস্পতিবার} mean? (\textquotedblleft{}Thursday\textquotedblright{}.) Say it once more.
\end{tcolorbox}
\section[\texorpdfstring{\bn{শুক্রবার} (shukrôbār)}{শুক্রবার (shukrôbār)}]{\bn{শুক্রবার} (shukrôbār) — Friday}

\label{lesson:BN-C91-shukrobar}



\begin{quote}
Before the new one: say the Bengali for Wednesday, then the Bengali for Thursday.
\end{quote}

\subsection*{You'll want to know: \bn{শুক্রবার}}
\textbf{\bn{শুক্রবার}} — \emph{shukrôbār} — \textquotedblleft{}Friday\textquotedblright{}.

\begin{cousinweb}
\textbf{\bn{শুক্রবার}} (\emph{shukrôbār}) is Sanskrit \emph{śukravāra}, \textquotedblleft{}Friday\textquotedblright{}, taken over whole as a learned loan, a tatsama. It is \emph{śukra}, \textquotedblleft{}Venus\textquotedblright{}, + \emph{vāra}, \textquotedblleft{}day of the week\textquotedblright{}.
\end{cousinweb}

\begin{grammarlens}[title={What you've built}]
One more word for this chapter.
\end{grammarlens}

\subsection*{Your turn}
\begin{itemize}
\raggedright
  \item \emph{Say it:} \emph{shukrôbār}
  \item \emph{Say it:} \emph{shukrôbār}, once more
  \item \emph{Say it:} say \emph{shukrôbār}
  \item \emph{Recall:} say the Bengali for Wednesday, then the Bengali for Thursday, then say \emph{shukrôbār} again
\end{itemize}

\begin{tcolorbox}[breakable,colback=teal!4,colframe=teal!35!black,title={Before you move on}]
What does \bn{শুক্রবার} mean? (\textquotedblleft{}Friday\textquotedblright{}.) Say it once more.
\end{tcolorbox}
